\begin{abstract}
Emerging work like BitNet~\cite{bitnet} is ushering in an era characterized by 1-bit Large Language Models (LLMs). In this study, we propose a new 1-bit LLM variant, designated as \textbf{\text{BitNet b1.58}}, where each parameter (weight) is constrained to take values from the ternary set \{-1, 0, 1\}. Remarkably, this model achieves parity with conventional full-precision (FP16 or BF16) Transformer LLMs of comparable size and trained on the same data in terms of perplexity and downstream performance. At the same time, it offers substantial improvements in latency, memory efficiency, throughput, and energy use. Importantly, this 1.58-bit LLM establishes an updated scaling law and methodology for developing subsequent generations of LLMs that prioritize both efficiency and effectiveness. Additionally, it provides a foundation for novel computational frameworks and inspires the design of dedicated hardware specialized for 1-bit LLM operations.
\end{abstract}

\section{The Era of 1-bit LLMs}

The landscape of AI has recently been transformed by the swift expansion in both scale and capability of Large Language Models (LLMs). While these systems have delivered outstanding outcomes across numerous natural language processing benchmarks, their mounting parameter counts have introduced obstacles to real-world deployment and heightened worries about energy usage and environmental costs.

To counteract these issues, one prevalent strategy involves post-training quantization, in which model parameters are reduced to lower bitwidths for inference~\cite{smoothquant, gptq, quip, quip_sharp}. This method slashes the storage and computation demands of LLMs by diminishing the precision of their weights and activations, and has driven a progression from 16-bit representations down to more compact forms, such as 4 bits~\cite{gptq, awq}. Despite their widespread adoption in production-grade models, these post-training quantization approaches are inherently limited and do not realize optimal performance.

Recent advancements in 1-bit model design, exemplified by BitNet~\cite{bitnet}, suggest a viable path for making LLMs more cost-effective without sacrificing capability. Traditional LLMs typically utilize 16-bit floating-point representations (either FP16 or BF16), with matrix multiplication constituting the majority of computational workload. As such, most computation overhead stems from floating-point operations, namely additions and multiplications. BitNet’s approach, however, shifts matrix multiplication to integer-only addition, thereby drastically lowering energy consumption for LLM workloads. Since many hardware platforms are constrained by power, these energy reductions can lead directly to improved computational speed.

Besides computation, the movement of model weights from DRAM to on-chip accelerator memory (like SRAM) can also become a costly bottleneck at inference time. Some approaches attempt to increase SRAM size for better throughput, though this is much more expensive than DRAM. 1-bit LLMs, compared to their full-precision counterparts, benefit from dramatically reduced memory requirements in both capacity and bandwidth, thus decreasing the time and expense of loading model weights and enabling quicker and more resource-efficient inference.

We present here a notable 1-bit LLM adaptation, \textbf{\text{BitNet b1.58}}, in which every model parameter is ternary and drawn from the set \{-1, 0, 1\}. By introducing 0 as an additional option to the originally binary 1-bit BitNet, we achieve a 1.58-bit system. \text{BitNet b1.58} preserves all principal advantages of BitNet, such as its nearly multiplication-free matrix computations that can be aggressively optimized, and maintains the same energy profile while delivering marked gains in memory efficiency, throughput, and latency relative to standard FP16-based LLMs. Importantly, \text{BitNet b1.58} offers two further benefits: firstly, it exhibits enhanced expressive power due to its explicit capacity for feature suppression—the 0 value enables effective filtering, leading to notable improvements in performance over strictly 1-bit LLMs. Secondly, our experimental findings indicate that, beginning at the 3B parameter scale, \text{BitNet b1.58} can achieve parity with full-precision (FP16) models in both perplexity and downstream task metrics when trained under identical configurations (e.g., model size, token count).

\section{BitNet b1.58}

\text{BitNet b1.58} derives from the BitNet family, employing a Transformer framework in which \emph{nn.Linear} layers are substituted with \emph{BitLinear} layers. The model is initialized with random weights and trained end-to-end, utilizing weights quantized to 1.58 bits and activations quantized to 8 bits. Several updates distinguish this version from the original BitNet, and these changes are outlined below.

\paragraph{Quantization Function.} For weight quantization to the ternary set \{-1, 0, +1\}, we utilize an \emph{absmean} quantization scheme. This method normalizes the weight matrix by dividing it by its mean absolute value, after which each element is discretized to the closest value in \{-1, 0, +1\}:

\begin{equation}
    \widetilde{W} = \mathrm{RoundClip}(\frac{W}{\gamma + \epsilon}, -1, 1),
\end{equation}

\begin{equation}
    \mathrm{RoundClip}(x, a, b) = \max(a, \min(b, \mathrm{round}(x))),
\end{equation}

\begin{equation}
    \gamma = \frac{1}{nm}\sum_{ij} |W_{ij}|.
\end{equation}

The approach for quantizing activations is largely consistent with the mechanism employed in BitNet. However, unlike prior practice, we refrain from rescaling activations before nonlinear operations to the interval $[0, Q_b]$. Instead, activation values are adjusted to fall within $[-Q_b, Q_b]$ on a per-token basis, thereby avoiding zero-point quantization. This adjustment streamlines both code implementation and system-level optimizations, and, as our results indicate, yields little to no negative impact on model accuracy.

\paragraph{LLaMA-alike Components.}

Given that the LLaMA~\cite{llama, llama2} architecture is widely adopted as the foundational structure for open-source large language models, we have structured \text{BitNet b1.58} to include several LLaMA-aligned elements. In particular, the design incorporates RMSNorm~\cite{rmsnorm}, SwiGLU~\cite{swiglu}, rotary positional encodings~\cite{rope}, and omits all bias parameters. This LLaMA-compatible design enables seamless integration with widely used open-source libraries such as Huggingface, vLLM~\cite{vllm}, and llama.cpp\footnote{\url{https://github.com/ggerganov/llama.cpp}} with minimal modifications required.

\section{Results}

We benchmarked \text{BitNet b1.58} against our implementation of FP16 \text{LLaMA LLM} across several model sizes. For consistency in the evaluation, all models were pre-trained for 100 billion tokens using the RedPajama dataset~\cite{redpajama}. Their zero-shot capabilities were assessed on a diverse set of language understanding tasks, such as ARC-Easy~\cite{arc}, ARC-Challenge~\cite{arc}, Hellaswag~\cite{hellaswag}, Winogrande~\cite{winoGrande}, PIQA~\cite{piqa}, OpenbookQA~\cite{openbookqa}, and BoolQ~\cite{boolq}. Additionally, validation perplexity was reported on WikiText2~\cite{wikitext2} and C4~\cite{c4}.

We also examined GPU runtime memory usage and inference latency for both \text{LLaMA LLM} and \text{BitNet b1.58}. These measurements utilized the FasterTransformer\footnote{\url{https://github.com/NVIDIA/FasterTransformer}} toolkit, which provides highly optimized support for large language model inference on GPUs. For \text{BitNet b1.58}, the 2-bit kernel from Ladder~\cite{ladder} was integrated. The primary metric reported was time per output token, as this constitutes the main inference expense.

Table~\ref{tab:ppl} presents the perplexity metrics and associated computational costs for \text{BitNet b1.58} in comparison with \text{LLaMA LLM}. The data indicate that starting at the 3B parameter scale, \text{BitNet b1.58} attains perplexity parity with its full-precision counterpart, while being 2.71 times faster and consuming just 28\% of the GPU memory (3.55$\times$ less). Notably, \text{BitNet b1.58} at 3.9B parameters surpasses the performance of \text{LLaMA LLM} 3B, with a 2.4$\times$ speedup and a 3.32$\times$ reduction in memory consumption.

Table~\ref{tab:zero-shot} details the zero-shot accuracy across end tasks. Evaluations adhered to the standard methodology of \emph{lm-evaluation-harness}\footnote{\url{https://github.com/EleutherAI/lm-evaluation-harness}}. As model sizes increase, the performance difference between \text{BitNet b1.58} and \text{LLaMA LLM} diminishes, with \text{BitNet b1.58} matching the accuracy of the full precision baseline at the 3B scale. In line with perplexity findings, the end-task results show that \text{BitNet b1.58} 3.9B exceeds the performance of \text{LLaMA LLM} 3B, while incurring much lower latency and memory costs. This evidences \text{BitNet b1.58} as a Pareto improvement relative to existing large language model approaches.

\paragraph{Memory and Latency}

We expanded our experiments by increasing the model sizes to 7B, 13B, and 70B, and subsequently assessed their associated costs. As depicted in Figure~\ref{fig:latency-memory-energy}, both latency and memory usage exhibit scaling behavior: performance gains become more pronounced as the model size increases. Notably, at 70B parameters, \text{BitNet b1.58} demonstrates a speed improvement of 4.1$\times$ over the \text{LLaMA LLM} reference model. This efficiency is attributable to the fact that the computational demand for \emph{nn.Linear} layers escalates with model size. In terms of memory, a parallel pattern emerges; since the embedding layer maintains full precision and its memory footprint constitutes a smaller fraction as models grow larger. Both latency and memory measurements employed a 2-bit kernel, indicating that additional optimizations could further lower the resource requirements.

\paragraph{Energy}

We performed an analysis of arithmetic operation energy consumption for both \text{BitNet b1.58} and \text{LLaMA LLM}, with emphasis on matrix multiplication as it represents the most significant cost in LLMs. Figure~\ref{fig:energy} breaks down the energy expenditure, showing that most operations in \text{BitNet b1.58} involve INT8 additions, whereas \text{LLaMA LLM} relies heavily on both FP16 addition and FP16 multiplication. Leveraging the energy model detailed in~\cite{energycost, pokebnn}, we find that \text{BitNet b1.58} achieves a 71.4-fold reduction in arithmetic energy consumption for matrix multiplication tasks on 7nm hardware. We also provide an assessment of the total end-to-end energy usage when processing sequences of 512 tokens. These findings indicate that, with growing model size, \text{BitNet b1.58} exhibits substantially greater energy efficiency relative to the FP16 \text{LLaMA LLM} baseline. This effect stems from the increasing dominance of \emph{nn.Linear} operations in larger models, with the relative impact of other system components diminishing accordingly.

\paragraph{Throughput}

To evaluate throughput, we conducted a comparison between \text{BitNet b1.58} and the \text{LLaMA LLM} model containing 70B parameters, utilizing two 80GB A100 GPUs. Pipeline parallelism~\cite{gpipe} was applied to make it feasible to deploy the 70B \text{LLaMA LLM} on this hardware. We progressively increased the batch size until we encountered GPU memory constraints, maintaining a sequence length of 512 throughout. The data in Table~\ref{tab:throughput} demonstrates that \text{BitNet b1.58} 70B accommodates a batch size up to 11 times larger than that of the \text{LLaMA LLM}, leading to an 8.9-fold improvement in throughput.

\textbf{\text{BitNet b1.58} enables a distinct scaling relationship linking model capability and inference expenditure}. Based on Figure~\ref{fig:latency-memory-energy} and \ref{fig:energy}, it is possible to equate certain model sizes in 1.58-bit precision with larger FP16 models in terms of efficiency metrics:

\begin{itemize}
    \item The 13B \text{BitNet b1.58} outperforms the 3B FP16 LLM with respect to latency, memory consumption, and energy efficiency.
    \item The 30B \text{BitNet b1.58} offers better latency, memory, and energy metrics than the 7B FP16 LLM.
    \item The 70B \text{BitNet b1.58} is more advantageous, considering latency, memory use, and energy needs, compared to the 13B FP16 LLM.
\end{itemize}

\paragraph{Training with 2T Tokens}

Training token count plays a pivotal role in the effectiveness of large language models. To assess token-level scaling properties of \text{BitNet b1.58}, we trained a \text{BitNet b1.58} model using 2T tokens, adhering to the data preparation approach of StableLM-3B~\cite{StableLM-3B-4E1T}, recognized as a leading open-source 3B parameter model. The evaluation encompassed a suite of tasks—Winogrande~\cite{winoGrande}, PIQA~\cite{piqa}, SciQ~\cite{sciq}, LAMBADA~\cite{lambada}, and ARC-easy~\cite{arc}—reporting zero-shot accuracy results in Table~\ref{tab:2t}. Where both accuracy and normalized accuracy metrics were reported, their mean value was used. StableLM 3b’s 2T token results are sourced from its official report. Our experimental evidence establishes that \text{BitNet b1.58} consistently outperforms its counterpart on all evaluated tasks, underscoring the strong generalization capability of 1.58-bit LLMs.

\section{Discussion and Future Work}

\textbf{1-bit Mixture-of-Experts (MoE) LLMs}

Mixture-of-Experts (MoE) architectures have emerged as a resource-efficient solution for large language models (LLMs), notably reducing computational FLOPs. Nevertheless, their substantial memory requirements and significant inter-chip communication costs continue to pose barriers to widespread deployment. 1.58-bit LLMs present viable solutions to these issues. Their smaller memory footprint directly translates to a lower number of hardware devices needed for MoE deployment. Furthermore, they notably diminish the expense associated with inter-device activation transfers. If all model parameters can reside within the memory of a single chip, the communication overhead may be effectively eliminated.

\textbf{Native Support of Long Sequence in LLMs}

Handling extended sequences natively within LLMs has become a fundamental requirement. One of the primary obstacles in processing long sequences lies in the substantial memory usage, especially due to KV cache storage during inference. \text{BitNet b1.58} makes notable progress in this direction by compressing activation precision from 16 bits down to 8 bits, thereby enabling double the context length for a given memory budget. Looking ahead, it should be feasible to losslessly reduce precision further, down to 4 bits or less for 1.58-bit LLMs, a direction we suggest for future investigation.

\textbf{LLMs on Edge and Mobile}

Deploying 1.58-bit LLMs holds considerable promise for edge and mobile environments, where memory and computational capabilities are inherently limited, constraining the complexity and scale of models that can be used. The advantages of decreased memory usage and reduced power requirements of 1.58-bit LLMs facilitate the execution of LLMs in such settings, opening up novel application domains that previously faced technical barriers. This progress stands to significantly augment the intelligence and functionality of edge and mobile platforms. Additionally, the increased CPU-friendliness of 1.58-bit LLMs, which aligns well with processor architectures prevalent in edge and mobile hardware, ensures that models such as \text{BitNet b1.58} can run with high efficiency, further elevating their potential in resource-constrained contexts.

\textbf{New Hardware for 1-bit LLMs}

The emergence of specialized accelerators such as those demonstrated by Groq\footnote{\url{https://groq.com/}} highlights the ongoing advancements in developing tailored hardware (e.g., LPUs) for efficient LLM computation. Building on this momentum, we advocate for the design of dedicated hardware and systems specifically tuned for the computation patterns of 1-bit LLMs, as made feasible by the BitNet~\cite{bitnet} paradigm.